\subsection{The unrestricted joint purity bound}\label{sec:fqe-joint-purity}

\begin{theorem}[Joint purity compatibility; consequence of QK]
\label{thm:fqe-joint-purity}
Every pure four-qubit state satisfies
\[
\boxed{\qquad T+\sum_{j=2}^4p_j^2\le3P-\frac94.\qquad}
\tag{JV}\label{eq:fqe-joint-purity}
\]
Equality holds exactly for HS cut spectra, cluster cut spectra, or products
of two Bell pairs, up to ordering of the cuts and eigenvalues.
In the HS case the states belong to the two HS LU orbits.
\end{theorem}
\begin{proof}
This is an elementary consequence of the stronger quadratic variance bound
\eqref{eq:fqe-quadratic-variance-proved}, together with $F\ge0$ and
$p_j\in[1/4,1]$. Set
\[
U=\sum_jp_j^2,\qquad G=3P-\frac94-T-U,\qquad
Q=F^2+\frac92(P-1)-6\bigl(U-P^2/3\bigr)\ge0.
\]
Since $T=(17P-12-F)/12$, direct algebra gives the useful identity
\[
12G=F(1-2F)+2(P-1)(3-2P)+2Q.
\tag{JI}\label{eq:fqe-joint-implication}
\]
It proves $G\ge0$ whenever $1\le P\le3/2$ and $0\le F\le1/2$.

For the remaining cases, use the purity box \eqref{eq:fqe-purity-box} of Section~\ref{sec:sup-qk}, the
maximum of the convex function $U$ at fixed $P$ over $[1/4,1]^3$.
Using $12G=19P-15+F-12U$, the first piece, for
$1\le P\le3/2$ and $F\ge1/2$, yields
\[
12G\ge F-\frac12+(P-1)(19-12P)\ge0.
\]
For $3/2\le P\le9/4$ the second piece yields
\[
12G\ge F+\frac12(2P-3)(31-12P)\ge0,
\]
and for $9/4\le P\le3$ the last piece yields
\[
12G\ge F-12P^2+67P-87\ge F+3>0.
\]
The last inequality follows from concavity of the quadratic and its
values $3,6$ at the two endpoints.

Equality in the first case forces $Q=0$, so the equality classification of
Corollary~\ref{thm:fqe-quadratic-variance} gives HS, cluster or Bell-pair
spectra. In the $F\ge1/2$ case equality requires
$P=1$, $F=1/2$, $U=3/8$, again giving $Q=0$ and the cluster
spectra. In the second box range equality requires
$P=3/2$, $F=0$ and $U=9/8$; its purity triple is
$(1,1/4,1/4)$, and the state is a product of two Bell pairs.
The final range is strict. All three listed types attain equality.
\end{proof}
